\chapter{Parallelism and domain decomposition}\label{ch:parallelism}
\chapterstatus{written}{2026-11-14}

\fml{} is written for distributed memory with MPI and shared memory with OpenMP, and every module also compiles and
runs without either. This chapter describes the runtime that sets them up, how the box is decomposed over the tasks
(slabs by default, pencils and Peano--Hilbert ranges as alternatives) and the communication helpers that the modules
share.

\begin{modulebox}{Global}
\nsp{FML}: \file{Global/Global.h} (runtime, reductions, \code{assert_mpi}), \file{Global/Domain.h} (\cls{Domain},
\cls{SlabDomain}, \cls{PencilDomain}, \cls{PeanoHilbertDomain}), \file{Global/DomainBuffer.h} (\cls{DomainBuffer}),
\file{Global/PeriodicSlabExchange.h}, \file{Global/ChunkedExchange.h}, \file{Global/MPICounts.h},
\file{Global/SystemMemory.h}.\par
Documentation: \file{Global/MPICounts.md}. Tests: \file{tests/test_domain.cpp}, \file{tests/test_mpi_counts.cpp}.
\end{modulebox}
\begin{modulebox}{PencilDomain, PeanoHilbertDomain}
\nsp{FML::GRID}: \file{PencilDomain/PencilGrid.h} (\cls{PencilGrid}), \file{PencilDomain/CPencilGrid.h}
(\cls{CPencilGrid}).
\nsp{FML::DOMAIN_DECOMPOSITION::PEANO_HILBERT}: \file{PeanoHilbertDomain/HilbertCurve.h} (\cls{HilbertCurve}),
\file{PeanoHilbertDomain/DistributedGrid3D.h} (\cls{DistributedGrid3D}).\par
Documentation: \file{PencilDomain/README.md}, \file{PeanoHilbertDomain/README.md}. Tests: \file{PencilDomain/tests},
\file{PeanoHilbertDomain/tests} and the \file{tests/test_*_pencil.cpp} suite.
\end{modulebox}

\section{The runtime}\label{sec:par-runtime}

With the compile flags \code{USE_MPI}, \code{USE_OMP} and \code{USE_FFTW} the library sets itself up automatically: a
singleton constructed at program start initialises MPI (with thread support), FFTW (with its MPI and threads interfaces)
and the globals, and finalises them at exit. With \code{NO_AUTO_FML_SETUP} the user calls \code{MPI_Init}, may set
\code{FML::communicator} to a sub-communicator (all communication of the library goes through it; default
\code{MPI_COMM_WORLD}) and then calls \code{FML::init_fml()} (and \code{FML::GRID::init_fftw}).

The globals are \code{FML::ThisTask} and \code{FML::NTasks}, the number of threads, whether MPI and threads can be combined
with FFTW, the local x-range of the particle domain \code{xmin_domain}, \code{xmax_domain} and the domain itself
(\cref{sec:par-domains}). On Linux the CPU affinity masks are checked: \code{NTasksThreadsShareCores} counts the tasks
that have fewer cores than OpenMP threads, a common misconfiguration of \code{mpirun}.

\paragraph*{Collectives and errors.} Inplace reductions of one value or an array (\code{SumOverTasks}, \code{MinOverTasks},
\code{MaxOverTasks}, \code{SumArrayOverTasks}, \ldots; arithmetic types reduce with \code{MPI_Allreduce} in bounded
chunks, others are gathered and reduced locally) and gathers (\code{GatherFromTasks}). Errors go through
\code{assert_mpi(condition, message)}, which prints the message and calls \code{MPI_Abort} rather than \code{abort} so that
the other tasks do not deadlock. Many modules can instead throw (e.g.\ \code{SPLINE_THROW_ON_ERROR},
\code{ODESOLVER_THROW_ON_ERROR}) so that a caller such as an MCMC can recover. With \code{MEMORY_LOGGING} every large
allocation of the standard containers is logged (\cref{ch:io-utilities}).

\paragraph*{Threads.} Grid loops are parallelised over \emph{rows} (lines of cells along the last axis) rather than slices,
because a task can have fewer slices than threads but not fewer rows. Particle loops use OpenMP with atomic or coloured
updates where they write to shared grids (\cref{ch:mass-assignment}).

\section{Domains}\label{sec:par-domains}

The unit box $[0,1)^D$ is decomposed by a \cls{Domain}, a \code{std::variant} of three types whose ownership queries
(\code{owner(pos)}, \code{owner_near(task, pos)}, \code{contains(task, pos)}) are inline switches without virtual calls:
\begin{itemize}
\item \cls{SlabDomain}: task $i$ owns $x\in[x_{\min,i},x_{\max,i})$, the ranges partitioning $[0,1)$ in rank order (empty
      slabs allowed). The default, given by \code{xmin_domain} and \code{xmax_domain} of each task (equal slabs
      $[i/N_{\rm tasks},(i+1)/N_{\rm tasks})$ from \code{init_fml}). Grids, particles and nearly all algorithms support it.
      \code{owner} binary-searches the edges; \code{owner_near} checks the neighbours first, where most leaving particles go.
\item \cls{PencilDomain}: a $P_0\times P_1$ process grid in $(x,y)$; the task with coordinates $(i_x,i_y)$ is rank
      $i_xP_1+i_y$ and owns a box in $x$ and $y$ and all of $z$. The decomposition of \cls{PencilGrid}.
\item \cls{PeanoHilbertDomain} (3D): contiguous ranges of the Peano--Hilbert key, built load-balanced from the particles.
\end{itemize}
The global instance \code{FML::domain} is set with \code{FML::set_domain} (collectively); \cls{MPIParticles} takes it when
created, so particles live on the task of their grid cells (\cref{ch:particles}). Code that needs slabs asserts it
(\code{require_slab}).

\paragraph*{Slabs.} A slab grid and the particles share the x-ranges; each task stores extra slices to the left and right
(ghosts) for the stencils of mass assignment, interpolation and finite differences. FFTW-MPI's slab transforms require
$N_{\rm mesh}$ divisible by the number of tasks and can use at most $N_{\rm mesh}$ tasks (one slice each).

\paragraph*{Pencils.} \cls{PencilGrid} splits two axes, so up to $N_{\rm mesh}^2$ tasks can be used and any $N_{\rm mesh}$ works
(the process grid from \code{MPI_Dims_create}, the edges on cells as evenly as possible). In 3D the transform is r2c along
$z$, a transpose within the row of tasks and an FFT along $y$, a transpose within the column and an FFT along $x$ (FFTW's
guru interface; the transposes pack cache-blocked blocks and exchange them with \code{MPI_Alltoallv}). In Fourier space a
task has all $k_x$ and parts of $k_y$ and $k_z$. On 2 tasks the forward and inverse FFTs take 0.9--1.3 times the slab time
and a whole N-body run 1.05--1.2 times; pencils pay off when the number of tasks approaches $N_{\rm mesh}$ or does not divide
it. The library, \file{SimulationCode} and the snapshot analyser are written against a common grid interface
(\cref{sec:grids-interface}) and run on slabs or pencils:

\begin{cppcode}[A run on pencils]
FML::set_domain(FML::pencil_domain());          // or pencil_domain(px, py)
auto grid = FML::GRID::make_grid<FML::GRID::PencilGrid<3>>(Nmesh, n_left, n_right);
FML::PARTICLE::MPIParticles<Particle> part;     // takes FML::domain
\end{cppcode}

\paragraph*{Peano--Hilbert ranges.} \cls{HilbertCurve} maps integer coordinates to 64-bit Hilbert keys with Skilling's
transpose algorithm \cite{Skilling2004} ($D=2$--6; in 3D up to 21 bits per coordinate). It is the one key implementation
of the library, used by \cls{AdaptiveMPIGrid} (\cref{sec:mg-amr}) and the particle domain. \code{PeanoHilbertDomain::build}
reduces a distributed prefix histogram of the keys ($2^{b_h}$ counters, $b_h\le20$) and cuts it into approximately
equal-particle contiguous key ranges at bin edges; identical keys always have the same owner, and empty ranges are allowed.
\cls{DistributedGrid3D} stores the cells of each task's key range in key order (not a rectangular block) with halo lists for
face or full cube stencils and collective remote reads. This decomposition is not yet used by \file{SimulationCode}.

\section{Communication patterns}\label{sec:par-communication}

\begin{itemize}
\item \emph{Ghost slices} of grids: \code{communicate_boundaries} sends the edge slices to the neighbours (forwarded from
      tasks further away when there are more ghost slices than local slices); \code{add_contribution_from_extra_cells}
      adds the ghosts back after mass assignment. \cls{GhostedCells} gives the same for per-cell data of any type
      (\cref{sec:grids-mpigrid}).
\item \emph{Particle migration} (\cref{ch:particles}): to the owner of the new position, with wide counters and bounded
      message sizes.
\item \emph{Buffer exchanges}: \cls{DomainBuffer} sends copies of the particles within a distance of a task's box to the
      neighbours (2 for slabs, up to 8 for pencils, corners included), optionally shifted across the periodic box, and maps
      values forward to the copies and back to the originals. It is used by the friends-of-friends, Delaunay and other particle
      methods. \code{append_periodic_slab_ghosts} (and \code{append_periodic_domain_ghosts} for pencils) does the same for records, skipping peers whose box is beyond the buffer.
\item \emph{Chunked exchange}: \code{exchange_records_chunked} delivers records to their destination task without a flattened send
      buffer or a full receive buffer, in rounds of at most 8\,MiB by default, with all tasks executing the same number of
      rounds.
\end{itemize}

\paragraph*{MPI counts.} MPI counts and displacements are \code{int}, so transfers above $2^{31}$ elements or bytes overflow
silently in naive code. \file{MPICounts.h} provides checked addition and multiplication, layouts that are verified before
any payload is sent (throwing \code{std::overflow_error} collectively), and gathers of fixed-size records. Particle
migration exchanges its counters as \code{unsigned long long} and chunks the payload by a message cap (default
\code{INT_MAX}); the other exchanges (sheets, FoF, ORIGAMI, Delaunay, void catalogues, MCMC summaries, \ldots) were audited
for the same limits (\file{MPICounts.md}).

\paragraph*{Memory.} \file{SystemMemory.h} reads the resident set size and its peak per task and prints their range over
tasks; the allocation log of \code{MEMORY_LOGGING} is described in \cref{ch:io-utilities}.

\begin{validation}[Validation: domains and communication]
\file{tests/test_domain.cpp} on 1, 2 and 4 tasks: the global Slab and its ranges; \code{owner} and \code{owner_near} of
slabs (with empty slabs), pencils and Peano--Hilbert ranges against brute force; \cls{MPIParticles} on every domain type
(\code{communicate_particles}, distributed and replicated input, checkpoints). \file{tests/test_mpi_counts.cpp}: exact
\code{INT_MAX} boundaries, overflow of aggregate layouts, one-rank failures handled collectively, empty ranks and a real
64-bit count exchange above \code{INT_MAX}. The pencil suite (\file{tests/test_*_pencil.cpp}) checks that the grid-based
modules (assignment, interpolation, LPT, spectra, cosmic web, voids, SPH, multigrid, \ldots) give the same results on slabs
and pencils.
\end{validation}
